% !TeX root = ../Chinese Philosophy.tex
\section{先秦诸子}

\subsection{前诸子时代和礼乐文明}
这一部分所称的前诸子时代，指的是夏商西周的礼乐文明时代。

礼乐文明的核心是制礼作乐（史料见《礼记·明堂位》）。礼的内容包括四个主要部分：\textbf{礼书}（《周礼》《礼记》《仪礼》），\textbf{名物}（衣服，饮食，车马balabala），\textbf{制度}（官，田，军，学，刑，宗法，祭祀），\textbf{礼仪}（冠，婚，丧，祭）

\subsection{孔子}

Mind：典型的“诸子学”这个概念是在汉代被建构起来的。

\subsubsection{孔子与儒家}

时代背景：礼崩乐坏，变革时代和秩序问题。

\begin{fancybox}{例子}
    孔子之时，周室微而礼乐坏，《诗》《书》缺。 （《孔子世家》）
    
    八佾舞于庭，是可忍也，孰不可忍也？
\end{fancybox}

在这样的时代背景之下，儒家的选择是一种突破与转化。

\subsubsection{名与礼}

要理解名与正名，首先要理解“名”的概念。在孔子哲学的视角下，“名”确定了人伦生活中各种关系，身份及其所构成的秩序。“名不正则言不顺，言不顺则事不成。”代表了通过名的划分，深入到人伦秩序的基础因素。

同时，礼乐秩序也是由“名”构建而成的，讲求“名”与“实”的对应。从以下引用中可以看到“礼之本末”。

\begin{fancybox}{例子}
    礼云礼云，玉帛云乎哉？  （论语《阳货》）

    尔爱（吝惜）其羊，我爱其礼。  （论语《八佾》）
\end{fancybox}

上面的例子中，礼的本质是礼乐秩序，礼的“末”即礼的仪式等外在表现。强调外在言行对礼的规定，同时强调人内部对礼制度体系的认同。在政治上和实践上主张道德意识的培养和道德教化。

\begin{fancybox}{培养道德意识}
    道之以政，齐之以刑，民免而无耻。道之以德，齐之以礼，有耻且格。（《论语·为政》）
\end{fancybox}

接下来谈论中庸之道：礼的规范与权宜。不追求纯粹抽象，以固化的规则为约束的系统。规则只能执行到一定程度，在规则之外需要动用其他的理性能力。

\begin{fancybox}{中庸哲学}
    中庸之为德，其至矣乎，民鲜久矣。   （《论语·雍也》）

    不偏之谓中，不易之谓庸。    （宋儒解中庸）

    过犹不及。     （《论语·先进》）
\end{fancybox}

\subsubsection{仁与礼}

克己复礼，仁德指的是按照礼的规范和约束要求自身的行为。仁和己的关系本身影射了一体两面的仁礼关系。

仁爱之始：亲亲相隐及其危险。朴素的理解：学习被爱，学习去爱别人。爱人是在家人之间首先培养的。当然这会导致大量的伦理问题。

\begin{fancybox}{仁爱之开始}
    樊迟问仁，子曰：“爱人”     （《论语·颜渊》）

    孝悌也者，其为人之本与！      （《论语·学而》）

    叶公语孔子曰：“吾党有直躬者，其父攘羊，而子证之。”       （《论语·子路》）
\end{fancybox}

为了解决上述伦理困境，衍生出的第一种解决方式是仁爱的推扩。

\begin{fancybox}{道德的推广}
    泛爱众，而亲仁。行有余力，则以学文。

    己所不欲，勿施于人。
\end{fancybox}

当然，在仁爱的推广中同样存在问题。问题在于这里的体察和对他人的伦理学态度是从自己的角度出发思考的。问题在于我所欲不一定是人所欲。

\subsubsection{教化}

教化问题可以分为几个层面的教化，第一个教化为政治教化。儒家考虑物质生活的富足，在物质生活丰富之后的最高目标是道德教化。“君子之德风，小人之德草，草上之风，必偃”说明儒家的政治主张是通过示范引导教化，不主张使用暴力和强制的手段。

教化方面，儒家有一些遗留的问题。其中重要的有思考秩序问题的两条路径人性与天道。

\begin{fancybox}{模糊之处}
    子贡曰：“夫子之文章，可得而闻也；夫子之言性与天道，不可得而闻也”
\end{fancybox}

\subsection{墨子}

墨子其人：鲁班和墨子的经典典故：

\begin{fancybox}{公输般}
    盖墨翟，宋之大夫，善守御，为节用。或曰并孔子时，或曰在其后。（《史记·孟子荀卿列传》）

    公输盘为楚造云梯之械，成，将以攻宋。子墨子闻之，起于齐，行十日十夜而至于郢，见公输盘。……子墨子解带为城，以牒为械，公输盘九设攻城之机变，子墨子九距之。公输盘之攻械尽，子墨子之守圉有馀。公输盘诎，而曰：“吾知所以距子矣，吾不言。”子墨子亦曰：“吾知子之所以距我，吾不言。”楚王问其故，子墨子曰：“公输子之意，不过欲杀臣。杀臣，宋莫能守，可攻也。然臣之弟子禽滑厘等三百人，已持臣守圉之器，在宋城上而待楚寇矣。虽杀臣，不能绝也。”楚王曰：“善哉！吾请无攻宋矣。”子墨子归，过宋。天雨，庇其闾中，守闾者不内也。（《墨子·公输》）
\end{fancybox}

\subsubsection{墨家思想}

墨家和儒家的关系较为复杂，在史料考证之外可以得到以下结论：墨家和儒家几乎同一时期，面对着相同的问题。墨家对于传统的态度应当为积极了解，在当时的社会背景之下取得了广泛的传播和认同。

\paragraph{天下之利}

追求天下人共同的物质财富和政治秩序。

\begin{fancybox}{天下之利}
子墨子言曰：仁人之所以为事者，必兴天下之利，除去天下之害，以此为事者也。（《墨子·兼爱中》）

子墨子言曰：今者王公大人为政于国家者，皆欲国家之富，人民之众，刑政之治。 （《墨子·尚贤上》）    
\end{fancybox}

墨家为了实现政治理想，提出了十大主张，在大量文献中对政治理想进行论述：

\textbf{兼爱、非攻、非乐、非命、节用、节葬、尚贤、尚同、天志、明鬼}

\paragraph{言立三表}

此问题与辩论相关，为了解决自身观点与他人观点之间可能发生的冲突，在学术和主张的竞争中坚持自身的政治理想，获得君主的认可。“言立三表”指的是三个要求：

\begin{itemize}
    \item 有本：要有古代圣王施政的原本和历史经验的依据。
    \item 有原：讲求体察民情，关注社会现实生活。
    \item 有用：具体进行了施政之后，政治主张在国家人民生活中体现出来。
\end{itemize}

这些问题是为了提供较为根本的原则和依据，在这条准则当中没有对逻辑学的探讨。

\begin{fancybox}{言立三表}
    言而毋仪，譬犹运钧之上而立朝夕者也，是非利害之辨，不可得而明知也。故言必有三表。何谓三表？子墨子言曰：有本之者，有原之者，有用之者。于何本之？上本之于古者圣王之事。于何原之？下原察百姓耳目之实。于何用之？废以为刑政，观其中国家百姓人民之利。此所谓言有三表也。（《墨子·非命上》） 
\end{fancybox}

\subsubsection{不相爱与兼爱}

\paragraph{不相爱}

对于社会的大量问题，墨子给出了简单的概括，将其概括为由于对私利的追求，导致人与人之间互相争夺，不互相爱护；或者爱是不够平等的私爱，只爱自己和国家。

\begin{fancybox}{材料}
    子墨子言曰：仁人之所以为事者，必兴天下之利，除去天下之害，以此为事者也。然则天下之利何也？天下之害何也？子墨子言曰：今若国之与国之相攻，家之与家之相篡，人之与人之相贼，君臣不惠忠，父子不慈孝，兄弟不和调，此则天下之害也。（《墨子·兼爱中》）
\end{fancybox}

墨家认为，害之生，以及社会问题来源于不相爱，或者爱有等差，自私自利。

\begin{fancybox}{材料}
    然则察此害亦何用生哉？以不相爱生邪？子墨子言：以不相爱生。今诸侯独知爱其国，不爱人之国，是以不惮举其国，以攻人之国。今家主独知爱其家，而不爱人之家，是以不惮举其家，以篡人之家。今人独知爱其身，不爱人之身，是以不惮举其身，以贼人之身。是故诸侯不相爱，则必野战。家主不相爱，则必相篡。人与人不相爱，则必相贼。君臣不相爱，则不惠忠。父子不相爱，则不慈孝。兄弟不相爱，则不和调。天下之人皆不相爱，强必执弱，富必侮贫，贵必敖贱，诈必欺愚。凡天下祸篡怨恨，其所以起者，以不相爱生也。是以仁者非之。（《墨子·兼爱中》）    
\end{fancybox}

\paragraph{兼爱}

墨子针对自己对于社会问题的论述，提出要用兼爱的方法加以解决，推崇博爱的作风。注意，墨子所谓“爱”和利益是紧密挂钩的，和日常语言中对于爱的理解不同。

\begin{fancybox}{材料}
    既以非之，何以易之？子墨子言曰：以兼相爱、交相利之法易之。然则兼相爱、交相利之法将奈何哉？子墨子言：视人之国，若视其国。视人之家，若视其家。视人之身，若视其身。是故诸侯相爱，则不野战。家主相爱，则不相篡。人与人相爱，则不相贼。君臣相爱，则惠忠。父子相爱，则慈孝。兄弟相爱，则和调。天下之人皆相爱，强不执弱，众不劫寡，富不侮贫，贵不敖贱，诈不欺愚。凡天下祸篡怨恨，可使毋起者，以相爱生也。是以仁者誉之。（《墨子·兼爱中》）
\end{fancybox}

\subsubsection{交相非与尚同}

墨子面对社会秩序的混乱，提出交相非的解释。认为所有处于自然状态的人，面对同一事情会有不同的评判标准，并认为自身的标准具备普世性。从而引发了多元，差异和对立（相亏害）。类似于西方政治哲学中的假设“一切人对一切人的战争”。

\begin{fancybox}{交相非}
    子墨子言曰：古者民始生、未有刑政之时，盖其语，人异义。是以一人则一义，二人则二义，十人则十义。其人兹衆，其所谓义者亦兹衆。是以人是其义，以非人之义，故交相非也。是以内者父子、兄弟作怨恶，离散不能相和合。天下之百姓，皆以水火、毒药相亏害。至有余力，不能以相劳。腐朽余财，不以相分。隐匿良道，不以相教。天下之乱，若禽兽然。（《墨子·尚同上》）
\end{fancybox}

解决社会问题，采用“尚同”的观点，统一标准，建构分级的政治体制（这其实与西周时期类似）。为了统一思想，对统一目的的描述就需要较为精准的表述，基于此，墨家发现了早期的逻辑学。尚同体制的原则构想，不同于人治，认为社会秩序的原则来源于所谓“天”。

\begin{fancybox}{尚同}
    夫明乎天下之所以乱者，生于无政长，是故选天下之贤可者，立以为天子。天子立，以其力为未足，又选择天下之贤可者，置立之以为三公......上以此为赏罚，甚明察以审信。
\end{fancybox}

墨子认为这种政治理想在公民当中的实践应该是通过赏罚完成的。而建构和维持政治制度，需要靠“尚贤”。儒家建构的政治制度，有血缘关系的考量，也有贤德的考量，概括而言就是亲亲+尊贤，从而设计了禅让制度。墨家则只有尊贤。所以两家都希望建立人与人之间相互关爱的关系，根本区别在于儒家主张从亲亲开始培养（爱有差等），“泛爱众”是一个最高的目标，但是仍然包括亲疏远近。墨家反对，将尚同作为通过赏罚推行的先验原则，没有差等和培养的过程。

\subsubsection{对儒家的批评}

\paragraph{节葬}

儒家认为应当厚葬久丧，墨家强烈反对，认为不应当花大量精力维护“亲亲”

\begin{fancybox}{节俭}
    今唯无以厚葬久丧者为政。国家必贫，人民必寡，刑政必乱。（《墨子·节葬下》）
\end{fancybox}

\paragraph{非乐}

墨家反对物质享受和不能带来物质财富的礼乐等精神文化活动，构成了对于儒家所继承的周代礼乐的主张的根本批评。因此对于传统墨子同样主张进行较大的改变。

\begin{fancybox}{非乐}
    子墨子言曰：仁之事者，必务求兴天下之利，除天下之害。将以为法乎天下，利人乎即为，不利人乎即止。且夫仁者之为天下度也，非为其目之所美，耳之所乐，口之所甘，身体之所安，以此亏夺民衣食之财，仁者弗为也。（《墨子·非乐上》）
\end{fancybox}

\subsection{老子}

老子其人有非常重的神秘感，《道德经》也有马王堆帛书，汉简，荆门郭店等等多种分章节或者不分章节的版本，多个版本之间还多有不同甚至矛盾之处。其人其书在本课程不做考证。

\subsubsection{老子思想的时代背景}

和孔子的时代背景是一样的，老子思想的时代背景也是社会大变革和随之而来的治理危机，具体而言就是周朝的统治和文化走向衰落。

\begin{fancybox}{周之衰：社会变革}
    朝甚除，田甚芜，仓甚虚，服文彩，带利剑，厌饮食，财货有余，是谓夸盗，非道也哉！

    民之饥，以其上食税之多，是以饥；民之难治，以其上之有为，是以难治；民之轻死，以其求生之厚，是以轻死。
\end{fancybox}

在社会变革时代需要对价值判断和价值体系进行重估，这也是老子思想的整体背景；与儒家主张的仁义孝悌和礼法规范有根本差别：

\begin{fancybox}{价值重估}
    大道废，有仁义；（智慧出，有大伪；）六亲不和，有孝慈；国家昏乱，有忠臣。

    绝圣弃智，民利百倍；绝仁弃义，民复孝慈；绝巧弃利，盗贼无有。（郭店本：绝智弃辩、绝伪弃诈、绝巧弃利）
\end{fancybox}

\subsubsection{道与德}

前诸子时代的思想核心为“德”，即人行为的准则。“道”是万物运行的规律，老子认为事物和人的运行规律和准则未必需要分离开来考虑。老子“道”的概念是从“德”中演化和发展而来的，并无明确的界限。在某些论证中，道的地位高于德，是更为根源性的准则。

\begin{fancybox}{道和德的关系}
    明道若昧，进道若退，夷道若类。上德若谷，太白若辱，广德若不足，建德若偷，质真若渝。大方无隅，大器晚成，大音希声，大象无形。

    道生之，德畜之。物形之，势（器）成之。是以万物莫不尊道而贵德……故道生之，（德）畜之，长之，育之，成之，熟之，养之，覆之。
\end{fancybox}

\subsubsection{玄德}

玄德是“道”所呈现出的运行规律的幽微而不居功的德性，所谓的“玄”指的是语言的有限性，虽然道和德是可以理解的，但是无法被语言清晰地表述出来。先分析以下对于玄德的论述：

\begin{fancybox}{玄德}
    生而不有，为而不恃，长而不宰，是谓玄德。
    
    大道泛兮，其可左右。万物恃之而生而不辞，功成而不名有。爱养万物而不为主，常无欲，可名于小。万物归焉而不为主，可名为大。是以圣人终不为大，故能成其大。
\end{fancybox}

玄德，针对的是历史上关于天地，先祖的崇拜。老子指出德行应为生下而不占有，养育而不带有自己的意志，具有“暗”的特征。实际上所谓的明德也可以具有不刻意张扬的“玄”的特征。\textbf{强调}：拥有玄德不是本身不可张扬，而是主动不去张扬。

\begin{fancybox}{明德}
    《诗》云：“予怀明德，不大声以色。”子曰：“声色之于以化民，末也。”（《中庸》）
\end{fancybox}

对于玄德的认识常常会犯阴谋论错误，譬如“国之利器不可以示人”，指的不是阴谋权斗，而是治国的方法和提要本身没有办法张扬出来，因为是“玄德”。

\paragraph{反，不自}
这是老子所谓玄德的核心特征，类似于黑格尔辩证法的自我否定。

否定的表达在老子的表述中是非常常见的。老子的否定不是通过否定他物来肯定自己，一旦能够界定自己，则脱离了“玄德”的境界。老子的否定是用来建构玄德的具体行为准则。不把自己放在万物的中心，不自以为是，才能达到“玄德”的要求，即“不自之自”。譬如下面一段，在对自身的否定中呈现了“后”的特征，而不是要否定“先”

\begin{fancybox}{反，不自}
    天长地久。天地所以能长且久者，以其不自生，故能长生。是以圣人后其身而身先，外其身而身存。不以其无私邪，故能成其私。

    贵以贱为本，高以下为基。是以侯王自谓孤寡、不谷，此其以贱为本邪？非乎？（\textbf{此为老子民本思想的表达}）

    柔胜刚，弱胜强，天下莫不知，莫能行。是以圣人言：受国之垢，是为社稷主；受国不祥，是为天下王。正言若反。

    反者道之动，弱者道之用。
\end{fancybox}

\paragraph{柔弱，居下，不争}

这是“不自”实践的核心特征。

\begin{fancybox}{实践}
    《庄子‧天下》：以濡弱谦下为表，以空虚不毁万物为实。

    上善若水。水善利万物而不争，处众人之所恶，故几于道。
    江海所以能为百谷王者，以其善下之，故能为百谷王。是以圣人欲上人，以其言下之；欲先人，以其身后之。是以处上而人不重，处前而人不害。是以天下乐推而不厌。以其不争，故天下莫能与之争。

    天之道，其犹张弓乎。高者抑之，下者举之，有余者损之，不足者与之。天道损有余而补不足。
    功成名遂身退，天之道。

    坚强者死之徒，柔弱者生之徒。
\end{fancybox}

\subsubsection{物}

老子对世界的认识：事物是多样的，事物是变化的，构成了一个丰富和多样的世界。事物的性质是辩证统一的。

\begin{fancybox}{物}
    夫天下神器，不可为也。为者败之，执者失之。凡物或行或随，或噤或吹，或强或剉，或培或堕。是以圣人去甚，去奢，去泰。

    有无之相生，难易之相成，长短之相形，高下之相倾，音声之相和，前后之相随。是以圣人处无为之事，行不言之教。
\end{fancybox}

\subsubsection{自然}

最重要的无为思想就是“辅自然”。自然有朴素的意涵。

\begin{fancybox}{道法自然}
    以辅万物之自然，而不敢为。
    
    道大，天大，地大，王亦大。域中有四大，而王居其一焉。人法地，地法天，天法道，道法自然。
\end{fancybox}

“道”的无为表达就是经典的\textbf{道法自然}，道虽然是最高的精神，核心是“因任万物之自然”。落实到行动中，则是对待欲望不特地用“知”去引导。也不应当追求名利和物质享受，让欲望处于合理的不受扰动的状态。也不应当引诱人的欲望，不论是欲望的对象，还是欲望在人与人之间关系的展开，都不应当受到引诱。

\begin{fancybox}{行为准则}
    不尚贤，使民不争；不贵难得之货，使民不为盗；不见可欲，使心不乱。是以圣人之治，虚其心，实其腹，弱其志，强其骨，常使民无知无欲。
    
    五色令人目盲，五音令人耳聋，五味令人口爽，驰骋田猎令人心发狂，难得之货令人行妨。是以圣人为腹不为目，故去彼取此。
    
    圣人自知不自见，自爱不自贵，故去彼取此。
\end{fancybox}

行为准则即：道应该是能够自主的自身，应当是不被其他事物影响和操控的自身。（师：对于消费主义的批评，不是因为消费人才能“成其所是”）

\subsubsection{道生万物}

这本质上是自然哲学问题，不仅仅限于物理世界的探讨，还要和人的行为准则建立联系。这决定了对人行为准则的理解。

\begin{fancybox}{道生万物}
    无名，万物之始；有名，万物之母。
    
    天下有始，以为天下母。既得其母，以知其子。既知其子，复守其母。没身不殆。

    天下之物生于有，有生于无。

    道生一，一生二，二生三，三生万物。万物负阴而抱阳，冲气以为和。
       
    昔得一者，天得一以清，地得一以宁，神得一以灵，谷得一盈，侯王得一以为天下正。
\end{fancybox}

注解：“一”有完整，独立，自主的意涵，是对自我观念的理解。

\subsection{孟子}

孟子其人：孟轲，驺人，受业于子思的门人。孟子生活在诸侯以富国强兵、合纵连横为务的战国时代，却主张唐虞三代的德政，因此其政治理想往往不为君主所用。退而与弟子整理《诗》《书》，阐述孔子之意，形成《孟子》七篇。

\begin{fancybox}{孟子其人}
    孟轲，驺人也。受业子思之门人。道既通，游事齐宣王，宣王不能用。适梁，梁惠王不果所言，则见以为迂远而阔于事情。……天下方务于合从连衡，以攻伐为贤，而孟轲乃述唐、虞、三代之德，是以所如者不合。退而与万章之徒序《诗》《书》，述仲尼之意，作《孟子》七篇。（《史记·孟子荀卿列传》）
\end{fancybox}

\subsubsection{孔孟之间}

\paragraph{儒家后学与儒家的分化}

孔子去世之后，弟子散游诸侯，或参与政治，或从事教学，或隐而不仕。儒家思想在传播过程中逐渐形成不同的传承和取舍。孟子继承孔子之业，同时根据战国时代的问题加以发展。

\begin{fancybox}{儒家后学}
    自孔子卒后，七十子之徒散游诸侯，大者为师傅卿相，小者友教士大夫，或隐而不见。……后陵迟以至于始皇，天下并争于战国，儒术既绌焉，然齐鲁之间，学者独不废也。于威、宣之际，孟子、荀卿之列，咸遵夫子之业而润色之，以学显于当世。（《史记·儒林列传》）
\end{fancybox}

儒家的内部并非只有一种解释，各派都自称继承了孔子的真正思想。这带来了如何理解和判定孔子之道的问题。

\begin{fancybox}{儒分为八}
    世之显学，儒、墨也。儒之所至，孔丘也。……自孔子之死也，有子张之儒，有子思之儒，有颜氏之儒，有孟氏之儒，有漆雕氏之儒，有仲良氏之儒，有孙氏之儒，有乐正氏之儒。……故孔、墨之后，儒分为八，墨离为三，取舍相反不同，而皆自谓真孔、墨，孔、墨不可复生，将谁使定后世之学乎？（《韩非子·显学》）
\end{fancybox}

\paragraph{儒家与其他学派}

孟子之“好辩”，与各家学说竞争、儒家需要申明自身立场的时代背景有关。孟子把维护仁义和人伦秩序看作辩论的目的，尤其批评杨朱的“为我”和墨子的“兼爱”。

\begin{fancybox}{予岂好辩哉}
    公都子曰：“外人皆称夫子好辩，敢问何也？”孟子曰：“予岂好辩哉？予不得已也。……圣王不作，诸侯放恣，处士横议，杨朱、墨翟之言盈天下。天下之言不归杨，则归墨。杨氏为我，是无君也；墨氏兼爱，是无父也。无父无君，是禽兽也。……杨墨之道不息，孔子之道不著，是邪说诬民，充塞仁义也。……我亦欲正人心，息邪说，距诐行，放淫辞，以承三圣者，岂好辩哉？予不得已也。能言距杨墨者，圣人之徒也。”（《孟子·滕文公下》）
\end{fancybox}

注意：“无君”“无父”是孟子从儒家立场作出的批评。孟子认为“为我”会损害对公共人伦秩序的责任，而无差等的兼爱会削弱亲亲的特殊地位。这与前面儒墨之争中，是否以亲亲作为仁爱起点的问题相联系。

\subsubsection{性善}

\paragraph{礼与心}

孟子讨论礼，强调把握事情的根本与先后轻重。仁者虽要爱人，但应首先重视亲近贤者；礼的实践也应先把握重要的人伦责任，不能只注意细小的礼仪规定。这承接了孔子关于“礼之本末”的讨论，并进一步引向礼背后的道德之心。

\begin{fancybox}{知务}
    孟子曰：“知者无不知也，当务之为急；仁者无不爱也，急亲贤之为务。尧、舜之知而不遍物，急先务也；尧、舜之仁不遍爱人，急亲贤也。不能三年之丧，而缌、小功之察；放饭流歠，而问无齿决，是之谓不知务。”（《孟子·尽心上》）
\end{fancybox}

\paragraph{何谓性：人禽之辨}

性可以初步理解为天生所具，但问题在于：哪些天生的能力和倾向构成人之为人的特征？告子以“生之谓性”来定义性，孟子则追问，这样是否会使犬、牛和人的性失去区别。

\begin{fancybox}{生之谓性}
    告子曰：“生之谓性。”

    孟子曰：“生之谓性也，犹白之谓白与？”曰：“然。”“白羽之白也，犹白雪之白；白雪之白，犹白玉之白欤？”曰：“然。”“然则犬之性犹牛之性，牛之性犹人之性欤？”（《孟子·告子上》）

    人之所以异于禽兽者，几希。（《孟子·离娄下》）
\end{fancybox}

这里的人禽之辨，指向人所具有的道德能力。人与禽兽的差别虽然“几希”，却是理解人性时应当重视的部分。

\paragraph{何谓性：性命之辨}

人的感官欲望和道德能力都可以说是天生所具，但孟子在二者之间作出了不同的强调。口腹、声色等欲望能否得到满足，受到外在条件和命的限制；仁义礼智的实践虽也涉及命，却应强调自身所具的性，不能以命为借口放弃道德责任。

\begin{fancybox}{性命之辨}
    孟子曰：“口之于味也，目之于色也，耳之于声也，鼻之于臭也，四肢之于安佚也，性也，有命焉，君子不谓性也。仁之于父子也，义之于君臣也，礼之于宾主也，智之于贤者也，圣人之于天道也，命也，有性焉，君子不谓命也。”（《孟子·尽心下》）
\end{fancybox}

注解：“臭”指气味。这里不是否认感官欲望属于性，而是说明君子在修养中如何看待欲望的满足与道德的实践。

\paragraph{四端之心}

孟子以“乍见孺子将入于井”为例说明，人会自发产生惊惧和恻隐之心。这一反应并非为了结交孩子的父母、获取名誉或避免恶名，因而显示了人本来就具有的“不忍人之心”。

% 课件第13页误标为《尽心下》，此段原文出自《公孙丑上》。
% 核对文本：https://www.shidianguji.com/book/SK0547/chapter/1jurkwe9ptova
\begin{fancybox}{四端之心}
    孟子曰：“人皆有不忍人之心。……所以谓人皆有不忍人之心者，今人乍见孺子将入于井，皆有怵惕恻隐之心，非所以内交于孺子之父母也，非所以要誉于乡党朋友也，非恶其声而然也。由是观之，无恻隐之心，非人也；无羞恶之心，非人也；无辞让之心，非人也；无是非之心，非人也。恻隐之心，仁之端也；羞恶之心，义之端也；辞让之心，礼之端也；是非之心，智之端也。人之有是四端也，犹其有四体也。有是四端而自谓不能者，自贼者也；谓其君不能者，贼其君者也。”（《孟子·公孙丑上》）
\end{fancybox}

四端是仁义礼智的开端，具体对应为：

\begin{itemize}
    \item 恻隐之心：不忍他人受苦，仁之端。
    \item 羞恶之心：对自身的不善感到羞耻，对不义感到厌恶，义之端。
    \item 辞让之心：尊重他人，能够辞让，礼之端。
    \item 是非之心：能够判断是非，智之端。
\end{itemize}

\textbf{注意}：人皆有四端，并不等于每个人已经具备充分实现的德行。“端”说明人有可以发展为善的起点，仍然需要后天的扩充和修养。

\paragraph{良能良知}

良能是“不学而能”，良知是“不虑而知”。孟子以儿童爱亲、长而敬兄为例，说明亲亲和敬长具有自发的基础。修养要把这种能力推达天下，与儒家从亲亲开始、逐步推广仁爱的思路相联系。

\begin{fancybox}{良能良知}
    人之所不学而能者，其良能也；所不虑而知者，其良知也。孩提之童，无不知爱其亲者，及其长也，无不知敬其兄也。亲亲，仁也；敬长，义也。无他，达之天下也。（《孟子·尽心上》）
\end{fancybox}

\paragraph{人所同然}

孟子从同类具有共同的感官偏好，推论人心也有共同认可的理与义。圣人与普通人同类，圣人只是先把握了人心所共同认可的道理。因此，道德能力具有普遍性，不能把圣人的德行看成普通人完全无法达到的事情。

\begin{fancybox}{心之所同然}
    孟子曰：“……故凡同类者，举相似也，何独至于人而疑之？圣人与我同类者。……故曰：口之于味也，有同耆焉；耳之于声也，有同听焉；目之于色也，有同美焉。至于心，独无所同然乎？心之所同然者何也？谓理也，义也。圣人先得我心之所同然耳。故理义之悦我心，犹刍豢之悦我口。”（《孟子·告子上》）
\end{fancybox}

注解：“耆”同“嗜”；“刍豢”指供食用的牲畜。这里强调人心能够认可并喜悦于理义。

\paragraph{人性善恶}

讨论人性善恶，首先要说明“性”的含义，以及“善恶”究竟指天生的能力、发展的可能，还是已经表现出来的行为。课件列出了几种不同的观点：

\begin{itemize}
    \item 告子：性无善无不善，性本身没有善恶的方向。
    \item 或者认为：性可以为善，也可以为不善，人的表现随着教化和环境而改变。
    \item 或者认为：有人的性善，有人的性不善，人与人的天性不同。
    \item 孟子：性善，人本来具有共同的道德能力和向善的倾向。
\end{itemize}

\begin{fancybox}{人性善恶的不同解释}
    公都子曰：“告子曰：‘性无善无不善也。’或曰：‘性可以为善，可以为不善。是故文、武兴则民好善，幽、厉兴则民好暴。’或曰：‘有性善，有性不善。是故以尧为君而有象，以瞽瞍为父而有舜，以纣为兄之子且以为君而有微子启、王子比干。’今曰‘性善’，然则彼皆非欤？”（《孟子·告子上》）
\end{fancybox}

告子以水可以向东流、向西流，比喻人性没有善恶方向。孟子则指出，水虽无东西之分，却有向下流的倾向。外力可以使水向上流，但不能因此把向上流看作水本身的性质；同样，人在某些条件下为不善，不足以否定其性善。

\begin{fancybox}{性善如水之就下}
    告子曰：“性犹湍水也，决诸东方则东流，决诸西方则西流。人性之无分于善不善也，犹水之无分于东西也。”孟子曰：“水信无分于东西，无分于上下乎？人性之善也，犹水之就下也。人无有不善，水无有不下。今夫水，搏而跃之，可使过颡；激而行之，可使在山。是岂水之性哉？其势则然也。人之可使为不善，其性亦犹是也。”（《孟子·告子上》）
\end{fancybox}

孟子进一步以“情”和“才”说明性善。“情”在此可理解为人本性的实际情状，“才”指人本来具备的能力。仁义礼智并非由外部强加给人，而有其内在基础；现实中的善恶差异，与是否思考、保存并充分发展这种能力有关。

\begin{fancybox}{非由外铄我也}
    孟子曰：“乃若其情，则可以为善矣，乃所谓善也。若夫为不善，非才之罪也。恻隐之心，人皆有之；羞恶之心，人皆有之；恭敬之心，人皆有之；是非之心，人皆有之。恻隐之心，仁也；羞恶之心，义也；恭敬之心，礼也；是非之心，智也。仁义礼智，非由外铄我也，我固有之也，弗思耳矣。故曰：‘求则得之，舍则失之。’或相倍蓰而无算者，不能尽其才者也。”（《孟子·告子上》）
\end{fancybox}

性善论所肯定的是人本有的向善基础。因此，现实中存在恶行与“性善”的主张可以并存，而心的修养正是把这种基础发展为实际德行的过程。

\subsubsection{心的修养}

\paragraph{扩充四端}

四端像初燃的火、初出的泉，具有继续发展的可能。修养需要把一时一事中表现出来的道德之心扩充，使它能够贯穿更多的关系和行动。仁爱的推扩从亲近之人开始，进一步及于他人，最终可以成为治理天下的基础。

\begin{fancybox}{扩充与推恩}
    凡有四端于我者，知皆扩而充之矣，若火之始然，泉之始达。苟能充之，足以保四海；苟不充之，不足以事父母。（《孟子·公孙丑上》）

    故推恩足以保四海，不推恩无以保妻子。（《孟子·梁惠王上》）

    老吾老，以及人之老；幼吾幼，以及人之幼。（《孟子·梁惠王上》）
\end{fancybox}

注解：“妻子”指妻子和子女。推恩将对自己亲人的关爱推及他人的亲人，也保留了从亲亲开始的次序。

\paragraph{养心寡欲}

欲望过多容易使人忽略本有的道德之心，因此孟子主张寡欲，以保存和养护此心。这里的重点是减少欲望对心的遮蔽，使道德判断和实践能够持续。

\begin{fancybox}{养心寡欲}
    孟子曰：“养心莫善于寡欲。其为人也寡欲，虽有不存焉者，寡矣；其为人也多欲，虽有存焉者，寡矣。”（《孟子·尽心下》）
\end{fancybox}

\paragraph{心之思}

人同样具备心与耳目等能力，却可能形成不同的生活取向。孟子把心称为“大体”，把感官称为“小体”：耳目容易受到外物引诱，心则能够思考并把握理义。“先立乎其大者”，就是先确立心的主导地位，使感官欲望不能夺去道德判断。

\begin{fancybox}{大体与小体}
    公都子问曰：“钧是人也，或为大人，或为小人，何也？”孟子曰：“从其大体为大人，从其小体为小人。”曰：“钧是人也，或从其大体，或从其小体，何也？”曰：“耳目之官不思，而蔽于物。物交物，则引之而已矣。心之官则思，思则得之，不思则不得也。此天之所与我者。先立乎其大者，则其小者不能夺也。此为大人而已矣。”（《孟子·告子上》）
\end{fancybox}

这里的“大人”“小人”，着重指人是否以心的思考和道德要求主导自身。心之思与扩充四端、养心寡欲相联系，共同说明性善如何落实为修养实践。

\subsubsection{仁政}

\paragraph{仁心与仁政}

仁政是“不忍人之心”在政治生活中的展开。人人都有仁心的开端，统治者应当把这种心扩充为对百姓的关怀，并落实为施政。由性善到修养，再到仁政，个人的道德能力与政治秩序由此建立联系。

\begin{fancybox}{不忍人之政}
    人皆有不忍人之心。先王有不忍人之心，斯有不忍人之政矣。以不忍人之心，行不忍人之政，治天下可运之掌上。（《孟子·公孙丑上》）
\end{fancybox}

\paragraph{制民之产}

仁政首先要保障百姓稳定的生计。“恒产”是稳定的生活来源，“恒心”是稳定的道德之心。孟子认为，士可以在没有恒产时仍然守住恒心，但普通百姓的道德生活需要基本的物质条件。若使百姓陷入贫困，待其犯罪后再加以惩罚，就是“罔民”。

\begin{fancybox}{恒产与恒心}
    无恒产而有恒心者，惟士为能。若民，则无恒产，因无恒心。苟无恒心，放辟邪侈，无不为已。及陷于罪，然后从而刑之，是罔民也。焉有仁人在位罔民而可为也？是故明君制民之产，必使仰足以事父母，俯足以畜妻子，乐岁终身饱，凶年免于死亡。然后驱而之善，故民之从之也轻。……此惟救死而恐不赡，奚暇治礼义哉？王欲行之，则盍反其本矣！（《孟子·梁惠王上》）
\end{fancybox}

注解：“罔民”指使百姓陷入罪中，犹如设网陷害。“制民之产”要求百姓能够赡养父母、养育妻儿，在丰年吃饱，在灾年免于死亡。在此基础上再施行教化，才有条件让百姓学习礼义。

\begin{fancybox}{制民之产的具体措施}
    五亩之宅，树之以桑，五十者可以衣帛矣。鸡豚狗彘之畜，无失其时，七十者可以食肉矣。百亩之田，勿夺其时，八口之家可以无饥矣。谨庠序之教，申之以孝悌之义，颁白者不负戴于道路矣。老者衣帛食肉，黎民不饥不寒，然而不王者，未之有也。（《孟子·梁惠王上》）
\end{fancybox}

具体措施包括保障田宅和农业生产，不耽误农时，发展桑蚕和畜养，并通过学校申明孝悌之义。物质生活的保障与道德教化共同构成仁政的内容。

\paragraph{人伦教化}

仅有衣食和安逸的生活，还不足以实现人之为人的道德生活。孟子强调在人伦关系中进行教化，明确父子、君臣、夫妇、长幼、朋友各自应有的关系和责任。这与前面孔子主张的先富而后教相联系，也回应了性善论中的人禽之辨。

\begin{fancybox}{人伦教化}
    人之有道也，饱食、暖衣、逸居而无教，则近于禽兽。圣人有忧之，使契为司徒，教以人伦：父子有亲，君臣有义，夫妇有别，长幼有叙，朋友有信。（《孟子·滕文公上》）
\end{fancybox}
